% 第八刀 GT 探针（2026-09-15）：\chardef 作用域 + char_given 显示面语义。
% 双引擎对拍（本探针只用 plain 已有原语）：
%   ntex-dvi probe-chardef-scope.tex out.dvi 2>&1 | grep '\['
%   pdftex -interaction=nonstopmode probe-chardef-scope.tex
% pdfTeX 期望（GT 权威，两引擎一致）：
%   [G:\char"41]             \global\chardef 组内绑定不随 \group_end 回滚；
%                            meaning 的 char_given 臂 = \char"<十六进制>
%   [local-rolled-back]      组内局部 chardef 出组回滚（cs 回到未定义 ≡ \relax）
%   [m:\char"7B]             同上，123 → "7B
%   [if:no][ifx:no]          \if 不透视 char_given（cs token ≠ 字符 token）；
%                            \ifx 对 \relax 亦假
%   [num:120]                无空格数字后续 `\the` 就地续数（12*10+\count0；
%                            此位 \gnum 是临时 \relax 绑定 → "can't use
%                            \relax after \the" 报一次后按 0 续；
%                            裸内部量如 `\count0` 不续、back_input）
\begingroup
  \global\chardef\gcharA=65
\endgroup
\message{[G:\meaning\gcharA]}
\begingroup
  \chardef\gcharB=66
\endgroup
\expandafter\ifx\csname gcharB\endcsname\relax
  \message{[local-rolled-back]}%
\else
  \message{[local-survived]}%
\fi
\chardef\gcharC=123
\message{[m:\meaning\gcharC]}
\message{[if:\if\gcharC B yes\else no\fi]}
\message{[ifx:\ifx\gcharC\relax yes\else no\fi]}
\count0=0
\chardef\gnum=12\the\count0
\message{[num:\the\gnum]}
\end
